% Transcription — fonds Alexandre Grothendieck, shelfmark 161-3, batch 1 (pages 1-20).
% Established from the facsimile archives/batches/161-3/batch-01.pdf (a mirror of
% https://grothendieck.umontpellier.fr/161-3.pdf), per the skill
% transcribe-grothendieck. Pages 2, 13 and 16 are blank — absent.
% Pass: Opus 5 (claude-opus-5), 2026-08-16 — first pass, unchecked against the pages by a human.
\documentclass[11pt,a4paper]{article}
% Le préambule commun à toutes les transcriptions du fonds Grothendieck.
%
% Il tient en une idée : **l'incertitude se compose, elle ne se résout pas.**
% Une transcription qui lisse un mot douteux a détruit la seule chose pour
% laquelle on la faisait — un lecteur doit pouvoir savoir, à chaque endroit,
% ce qui était sur le feuillet et ce qui a été ajouté. Les six macros qui
% suivent sont donc l'appareil critique, et non de la décoration.
%
% Les mêmes commandes sont reconnues par `scripts/render.mjs`, qui produit la
% vue de lecture du site. Une macro ajoutée ici doit l'être aussi là-bas,
% faute de quoi le rendu échouera bruyamment — ce qui est voulu : un rendu
% silencieusement faux est pire qu'un rendu absent.

\ProvidesPackage{grothendieck}[2026/08/08 Transcriptions du fonds Grothendieck]

\usepackage{amsmath,amssymb,amsthm}
\usepackage{mathrsfs}
\usepackage{stmaryrd}
% Les diagrammes commutatifs sont partout dans ce fonds : c'est la langue
% même de la géométrie algébrique de Grothendieck, et une transcription qui
% les rendrait en prose ne transcrirait rien.
\usepackage{tikz-cd}
% Français par défaut — c'est la langue des feuillets — mais l'anglais est
% chargé aussi : la traduction et le résumé sont des documents anglais, et la
% typographie française (espaces avant les deux-points) y serait une faute.
% Ces documents basculent par \selectlanguage{english} en tête de corps.
\usepackage[english,french]{babel}
\usepackage{fontspec}
\usepackage{geometry}
\geometry{a4paper,margin=25mm,marginparwidth=28mm}
\usepackage{marginnote}
\usepackage{xcolor}
% ulem, not soul. `soul`'s \st and \ul are built on a character-by-character
% scan that XeLaTeX's font machinery breaks: the document fails with an
% undefined control sequence rather than degrading. ulem does the same two jobs
% and is Unicode-engine safe. `normalem` stops it hijacking \emph, which the
% transcriptions use for Grothendieck's own emphasis.
\usepackage[normalem]{ulem}
\usepackage{hyperref}
\hypersetup{colorlinks=true,linkcolor=black,urlcolor=black,pdfborder={0 0 0}}

% --- Métadonnées du lot -----------------------------------------------------
% Reprises telles quelles par le script de rendu, qui les affiche en tête de
% la vue de lecture. Elles fixent ce que le fichier prétend couvrir : sans
% elles, une transcription flotte sans point d'attache dans un fonds de seize
% mille feuillets.

\newcommand{\folder}[1]{\gdef\grothFolder{#1}}
\newcommand{\batch}[1]{\gdef\grothBatch{#1}}
\newcommand{\leaves}[2]{\gdef\grothFirst{#1}\gdef\grothLast{#2}}
\newcommand{\foldertitle}[1]{\gdef\grothFoldertitle{#1}}
\gdef\grothFolder{?}\gdef\grothBatch{?}\gdef\grothFirst{?}\gdef\grothLast{?}\gdef\grothFoldertitle{}

% --- Le résumé --------------------------------------------------------------
%
% Il ouvre la lecture modernisée, sous le titre « Résumé », et il est composé
% en petit. Ce n'est pas un ornement : c'est le seul passage du document qui
% ne soit pas des mathématiques, et un lecteur doit pouvoir le reconnaître
% avant de l'avoir lu — puis le sauter s'il connaît déjà le sujet, ou s'y
% arrêter s'il ne le connaît pas. Le filet qui le suit marque la reprise du
% corps.

\newenvironment{resume}
  {\small\setlength{\parskip}{0.45em}}
  {\par\medskip\hrule\medskip}

% --- Mots-clés ---------------------------------------------------------------
%
% En anglais, en clôture du résumé de la lecture modernisée : le vocabulaire
% moderne sous lequel chercher ce que les feuillets construisent. Le manifeste
% du site (`npm run manifest`) les extrait pour étiqueter la cote entière —
% cette ligne est la source unique des tags, il n'y a pas de fichier à côté.
\newcommand{\keywords}[1]{\par\smallskip\noindent
  {\footnotesize\textsc{Keywords} --- \textit{#1}}\par}

% --- L'appareil critique ----------------------------------------------------

% Le numéro de feuillet, en marge. C'est l'ancre qui permet de régler un
% désaccord sur une formule en regardant *un* feuillet plutôt qu'une chemise
% de six cents. Le site s'en sert aussi pour tourner les pages du fac-similé
% quand on fait défiler la transcription.
\newcommand{\leaf}[1]{%
  \marginnote{\footnotesize\color{gray!70}\textsf{#1}}%
  \label{leaf:#1}\ignorespaces}

% Illisible : on ne devine pas. Un mot inventé qui se lit comme les autres est
% le pire résultat possible.
\newcommand{\ill}{\textcolor{red!60!black}{[\,\ldots\,]}}

% Lecture douteuse : le mot est donné, et signalé comme incertain.
\newcommand{\uncertain}[1]{\uline{#1}}

% Ajout de l'éditeur — une lettre manquante, un mot sous-entendu.
\newcommand{\add}[1]{\textcolor{blue!50!black}{[#1]}}

% Biffé par Grothendieck lui-même. Ce qu'il a rayé fait partie du document :
% une rature dit souvent où la pensée a changé de direction.
\newcommand{\struck}[1]{\sout{#1}}

% Note du transcripteur, sur l'état matériel du feuillet.
\newcommand{\note}[1]{\par\smallskip{\small\color{gray!60!black}\textit{[#1]}}\par\smallskip}

% Note marginale de Grothendieck, distincte du corps du texte.
\newcommand{\marginal}[1]{\par\smallskip\noindent\hspace{1em}%
  {\small\color{gray!60!black}#1}\par\smallskip}

% --- Titre ------------------------------------------------------------------

\AtBeginDocument{%
  \begin{center}
    {\large\bfseries Fonds Alexandre Grothendieck --- cote n\textsuperscript{o}~\grothFolder}\\[2pt]
    {\itshape\grothFoldertitle}\\[4pt]
    {\small feuillets \grothFirst--\grothLast\ (lot \grothBatch)}\\[2pt]
    {\footnotesize\color{gray!60!black}Transcription établie d'après le fac-similé numérisé
     par l'Université de Montpellier.\\
     Les lectures incertaines sont soulignées, les illisibles notées [\,\ldots\,],
     les ajouts entre crochets.}
  \end{center}
  \vspace{1em}}

\endinput


\folder{161-3}
\batch{1}
\pages{1}{20}
\dating{[vers 1963-1973]}
\watermark{Édition de démonstration}
\foldertitle{Topos : notes manuscrites (s.d.)}

\begin{document}

\page{1}
Titre porté au crayon sur l'onglet de la chemise : « Topos ».
\note{chemise du dossier}

\section*{Produits de topos}

\page{3}
$\Pi(E,F) \simeq$ faisceaux sur $E$ à valeurs dans $F$, i.e. foncteurs
$E^{\circ} \to F$ \struck{\ill{}} aux $\varprojlim$, i.e. qui ont un adjoint à
gauche

$\simeq$ faisceaux sur $F$ à valeurs dans $E$, i.e. foncteurs $F^{\circ} \to E$
commutant aux $\varprojlim$, i.e. qui ont un adjoint à g.

$\simeq$ foncteurs $E^{\circ} \times F^{\circ} \to (\mathrm{Ens})$ qui
\struck{\ill{}} en \emph{chaque variable} aux $\varprojlim$ quelc.
\note{le mot biffé n'est remplacé par aucun autre}

\marginal{si $E \simeq \widetilde{C}$}
$\simeq$ faisceaux sur $C$ à valeurs dans $F$, i.e. foncteurs
$\varphi : C^{\circ} \to F$ tels que pour tout $R \to X$ couvrant dans $C$, on a
$\varphi(X) \xrightarrow{\sim} \varprojlim_{C/R} \varphi(Z)$, i.e. tels que
$\forall\, y \in F$, $X \mapsto \operatorname{Hom}_{F}(\varphi(X), y)$ \ill{}
sur $C$ est un faisceau, (i.e. \struck{\ill{}} foncteurs
$\varphi^{\circ} : C \to F^{\circ}$ qui sont « continus » [i.e. se prolongent en
foncteurs $\widetilde{C} \to F^{\circ}$ commutant aux $\varprojlim$ [quelc.]])

\marginal{si $F \simeq \widetilde{D}$}
$\simeq$ faisceaux sur $D$ à valeurs dans $E$, i.e. foncteurs
$\psi : D^{\circ} \to E$ tels que \ill{}, i.e. foncteurs « continus »
$\psi^{\circ} : D \to E^{\circ}$

$\simeq$ foncteurs $C^{\circ} \times D^{\circ} \to (\mathrm{Ens})$ qui sont des
« bifaisceaux » i.e. tels que ce soient des faisceaux séparément en chaque
variable.

\struck{\ill{} $E \to E'$, $v : F \to F'$ \ill{}}

\textbf{Prop. 1} \quad $\Pi(E,F)$ est un topos.

On peut supposer $E = \widetilde{C}$, $F = \widetilde{D}$, $C$ et $D$ petites
stables par $\varprojlim$ finies.

a) Cas des topologies discrètes : $\Pi(\widehat{C}, \widehat{D}) \simeq
\widehat{C \times D}$ : c'est un topos.

b) Plus généralement, $\Pi(\widehat{C}, \widetilde{D}) \simeq
\operatorname{Hom}(C^{\circ}, \widetilde{D})$ est un topos, \struck{de même}
$\Pi(\widetilde{C}, \widehat{D}) \simeq \operatorname{Hom}(D^{\circ},
\widetilde{C})$.

\struck{Et} L'inclusion naturelle $\Pi(\widehat{C}, \widetilde{D})
\hookrightarrow \Pi(\widehat{C}, \widehat{D})$ (qui à $\varphi$ associe
$i_{D} \circ \varphi$, où $i_{D} : \widetilde{D} \to \widehat{D}$ est
l'inclusion) est le foncteur image directe pour un morphisme de topos, dont le
foncteur image inverse $\Pi(\widehat{C}, \widehat{D}) \to \Pi(\widehat{C},
\widetilde{D})$ est donné par $\varphi \mapsto a_{D} \circ \varphi$, où
$a_{D} : \widehat{D} \to \widetilde{D}$ est l'adjoint à g. de $i_{D}$.

Donc $\Pi(\widehat{C}, \widetilde{D})$ est un sous-topos
(\uncertain{str. pleine}) dans $\Pi(\widehat{C}, \widehat{D})$, et par la raison
symétrique $\Pi(\widetilde{C}, \widehat{D})$ aussi. Mais $\Pi(\widetilde{C},
\widetilde{D})$ est l'intersection de ces deux, donc c'est un sous-topos de
$\Pi(\widehat{C}, \widehat{D})$

\page{4}
— a fortiori un topos.

\struck{\ill{}}

On va définir, pour deux morphismes de topos
\[
u : E \longrightarrow E', \qquad v : F \longrightarrow F'
\]
un morphisme de topos
\[
\Pi(u,v) : \Pi(E,F) \longrightarrow \Pi(E',F')
\]
dont l'interprétation intrinsèque symbolique est ainsi, en termes des
« bifaisceaux » $E^{\circ} \times F^{\circ} \to (\mathrm{Ens})$, et
$E'^{\circ} \times F'^{\circ} \to (\mathrm{Ens})$,

\struck{\ill{}}
\[
\Pi(u,v)_{*}(\varphi)(X',Y') = \varphi(u^{*}X', v^{*}Y')
\qquad : \Pi(E,F) \longrightarrow \Pi(E',F')
\]
\[
\Pi(u,v)^{*}(\varphi)(X,Y) \overset{?}{=}
\]

Pour prouver que \struck{\ill{}} le foncteur précédent est associé à un
morphisme de topos, nous sommes ramenés, en le factorisant en
\[
\Pi(E,F) \longrightarrow \Pi(E,F') \longrightarrow \Pi(E',F')
\]
au cas où l'un des deux morphismes $u$, $v$ est l'identité ; soit $u$, d'où à
étudier
\[
\Pi(\mathrm{id}_{E}, v)_{*} : \Pi(E,F) \longrightarrow \Pi(E,F')
\]
On peut supposer $v : F \to F'$ défini par un morphisme de sites
$v_{0}^{*} : C' \to C$, $C$, $C'$ stables par \struck{\ill{}} $\varprojlim$
finies, \uncertain{munis d'une} top. moins fine que la top. can. Donc on a
\struck{\ill{}} un diagramme \struck{commutatif} de topos

\struck{$F \xrightarrow{\ \widetilde{v}_{0}\ } F'$}

\begin{tikzcd}
\widetilde{C} = F \arrow[r, "\widetilde{v}_{0}"] \arrow[d, hook] & \widetilde{C}' = F' \arrow[d, hook] \\
\widehat{C} \arrow[r, "\widehat{v}_{0}"] & \widehat{C}'
\end{tikzcd}

et par suite un diagramme. $E = \widetilde{C}$, et on a alors
$\Pi(\mathrm{id}_{E}, v)_{*}$ s'inscrit dans un diagramme commutatif

\begin{tikzcd}[column sep=small, row sep=small, nodes={font=\scriptsize}]
\Pi(\widetilde{C}, F) \arrow[r, "{\Pi(\mathrm{id}_{\widetilde{C}}, v)_{*}}"] \arrow[d, "{\Pi(i_{C}, \mathrm{id}_{F})_{*}}"'] & \Pi(\widetilde{C}, F') \arrow[d, "{\Pi(i_{C}, \mathrm{id}_{F'})_{*}}"] \\
\Pi(\widehat{C}, F) \arrow[r, "{\Pi(\mathrm{id}_{\widehat{C}}, v)_{*}}"'] & \Pi(\widehat{C}, F')
\end{tikzcd}

\page{5}
Nous avons déjà \struck{\ill{}} dans le \uncertain{dessin} de (prop. 1) vérifié
que les flèches verticales sont des inclusions de sous-topos [inclusion d'un
sous-topos intersection dans un des facteurs] donc pour prouver que
$\Pi(\mathrm{id}_{\widetilde{C}}, v)_{*}$ est \struck{\ill{}} le foncteur
\uncertain{adjoint} pour un morphisme de topos, il suffit de vérifier qu'il en
est ainsi de $\Pi(\mathrm{id}_{\widehat{C}}, F')_{*}$. En effet,
\note{le second argument de $\Pi(-,-)$ est partout ailleurs un morphisme ; la
page porte bien $F'$}

On aura alors un morphisme de topos
\[
\Pi(\widetilde{C}, F) \longrightarrow \Pi(\widehat{C}, F') ,
\qquad\text{et}
\]
\note{une flèche oblique, sous la précédente, la fait passer par
$\Pi(\widehat{C}, F)$}

on voit (en vérifiant sur les foncteurs images directs) qu'il se factorise par
le sous-topos $\Pi(\widetilde{C}, F')$, le foncteur image directe associé au
morphisme de \ill{} topos obtenu $\Pi(\widetilde{C}, F) \to
\Pi(\widetilde{C}, F')$ étant $\Pi(\mathrm{id}_{\widetilde{C}}, v)_{*}$.
\note{une large tache d'encre couvre ici trois demi-lignes}

Or $\Pi(\mathrm{id}_{\widehat{C}}, \ill{})$ s'identifie au foncteur
\[
\operatorname{Hom}(C, F) \longrightarrow \operatorname{Hom}(C, F')
\]
défini par \struck{\ill{}} $\varphi \mapsto v_{*} \circ \varphi$ [car la
bijection associée à $\varphi$ est $\varphi(X,Y) \to
\operatorname{Hom}_{F}(Y, \varphi(X))$, qui se transforme par
$\operatorname{Hom}(X,Y') \to \operatorname{Hom}_{F}(v^{*}Y', \varphi(X))
\simeq \operatorname{Hom}_{F'}(Y', (v_{*}\varphi)(X))$, et pour
$\varphi' = v_{*}\varphi : C \to F'$, c'est bien la bijection sur
$C \times F'$ associée à $\varphi'$], dont l'adjoint à droite est
$\varphi \mapsto v^{*} \circ \varphi$, et celui-ci est bien exact à g. cqfd.
\note{c'est l'adjoint à gauche de $\varphi \mapsto v_{*}\circ\varphi$ que
demande la suite ; la page porte « droite »}

Prenons $v : F \to P$ (topos ponctuel $\simeq \widehat{e}$, $e$ catégorie
\uncertain{finale}) et notant que $\Pi(E,P) \cong E$, on déduit un morphisme
can.
\[
\Pi(E,F) \longrightarrow E
\]
et on trouve de même un morphisme can. de topos
\[
\Pi(E,F) \longrightarrow F
\]
d'où un morphisme de topos
\[
\Pi(E,F) \longrightarrow E \times_{\mathrm{top}} F
\]
évidemment 2-fonctoriel en $E$, $F$. On veut prouver que c'est une
\struck{\ill{}} équivalence, i.e. que $\Pi(E,F)$ « est » un 2-produit des

\page{6}
topos $E$, $F$. On suppose comme dans la dém. de prop. 1 $E = \widetilde{C}$,
$F = \widetilde{D}$, \struck{donc} on a vu que \ill{}
\[
\Pi(E,F) = \Pi(\widetilde{C}, \widetilde{D}) = \struck{\ill{}}\;
\Pi(\widetilde{C}, \widehat{D}) \cap \Pi(\widehat{C}, \widetilde{D})
\]
(intersection des sous-topos de $\Pi(\widehat{C}, \widehat{D}) \simeq
\widehat{C \times D}$)

et il est clair que les deux foncteurs $\Pi(\widetilde{C}, \widehat{D})$ et
$\Pi(\widehat{C}, \widetilde{D})$ sont les images inverses des sous-topos
$\widetilde{C} \subset \widehat{C}$ resp. $\widetilde{D} \subset \widehat{D}$
par les « projections » $\Pi(\widehat{C}, \widehat{D}) \to \widehat{C}$ resp.
$\Pi(\widehat{C}, \widehat{D}) \to \widehat{D}$. Cela nous ramène à
\uncertain{savoir} prouver que $\Pi(\widehat{C}, \widehat{D})$ est un
2-produit des topos $\widehat{C}$, $\widehat{D}$. Or pour un topos
\struck{$E$}$X$, on a
\begin{align}
\operatorname{Hom}_{\mathrm{top}}(X, \widehat{C}) &\cong \operatorname{Hom}'(C, X)^{\circ} \nonumber\\
\operatorname{Hom}_{\mathrm{top}}(X, \widehat{D}) &\cong \operatorname{Hom}'(D, X)^{\circ} \nonumber
\end{align}
où le $\operatorname{Hom}'$ désigne les foncteurs exacts à g. [NB \ill{}
suppose $C$ et $D$ stables par $\varprojlim$ finies], donc on a
\begin{align}
\operatorname{Hom}_{\mathrm{top}}(X, \widehat{C}) \times \operatorname{Hom}_{\mathrm{top}}(X, \widehat{D})
&\cong \bigl[\operatorname{Hom}'(C,X) \times \operatorname{Hom}'(D,X)\bigr]^{\circ} \nonumber\\
&\cong \bigl(\operatorname{Hom}'(C \times D, X)\bigr)^{\circ} \nonumber\\
&\cong \operatorname{Hom}_{\mathrm{top}}(X, \widehat{C \times D}) \nonumber
\end{align}
cqfd

[si $C$ et $D$ sont stables par $\varprojlim$ finies, \struck{il} $C$ en est de
même de $C \times D$ (Attention, \ill{} des variables, non pas \ill{} les
variables séparément) et les foncteurs exacts à g. $C \times D
\xrightarrow{\ \sim\ } X$ [$X$ une catégorie quelc.] correspondant aux
\struck{deux} couples de foncteurs exacts à gauche $C \xrightarrow{\ u\ } X$,
$D \xrightarrow{\ v\ } X$, et leurs associés \struck{\ill{}} $(a,b) \mapsto
u(a) \times v(b)$, et à $w$ les foncteurs $a \mapsto w(a, e_{D})$, $b \mapsto
w(e_{C}, b)$

\page{7}
a) $\widehat{C} \times \widehat{C}' \simeq \widehat{C \times C'}$

$C$, $C'$ petites catégories avec $\varprojlim$ finies

Plus généralement, si on a une famille de catégories $C_{i}$ petites avec
$\varprojlim$ finies, soit $\Pi' C_{i} \subset \Pi C_{i}$ la sous-catégorie
\struck{\ill{}} pleine produit formée des systèmes $(x_{i})_{i \in I}$
($x_{i} \in \operatorname{Ob} C_{i}$) tel que $x_{i} \simeq e_{C_{i}}$ pour
presque tout $i$, alors
\[
\prod_{i} \widehat{C_{i}} \;\approx\; \widehat{{\textstyle\prod}' C_{i}}
\]

\marginal{$\varphi, \varphi'$ exacts à gauche ; $C \xrightarrow{\varphi} X$,
$C' \xrightarrow{\varphi'} X$ ($X$ catégorie aux $\varprojlim$ finies) se
factorisent en $C \xrightarrow{i_{1}} C \times C' \xrightarrow{f} X$,
$C' \xrightarrow{i_{2}} C \times C'$, où $i_{1}(x) = (x, e_{C'})$,
$i_{2}(x') = (e_{C}, x')$, par $f(x,x') \simeq \varphi(x) \times \varphi'(x')$ ;
$f(\varprojlim(x_{i}, x'_{i})) = f(\varprojlim x_{i}, \varprojlim x'_{i}) =
\varphi(\varprojlim x_{i}) \times \varphi'(\varprojlim x'_{i}) =
\varprojlim_{i} \varphi(x_{i}) \times \varprojlim_{i} \varphi'(x'_{i}) =
\varprojlim_{i} (\varphi(x_{i}) \times \varphi'(x'_{i}))$ OK}

b) Soient $E_{i}$ des topos ($i \in I$), $E'_{i} \underset{\alpha_{i}}{\subset}
E_{i}$ sous-topos, \ill{} 2-produit strict

Supposons $E = \prod_{i}^{\mathrm{top}} E_{i} \xrightarrow{\ \mathrm{pr}_{i}\ }
E_{i}$ existe

Soit pour tout $i$, $F_{i} \subset E$ le sous-topos
$\mathrm{pr}_{i}^{-1}(E'_{i})$

Soit $E' = \bigcap F_{i}$ le sous-topos intersection

Donc on a $E' \xrightarrow{\ \mathrm{pr}'_{i}\ } E'_{i}$, $E' \simeq
\prod^{\mathrm{top}} E'_{i}$

Ceci fait de $E'$ un 2-produit des $E'_{i}$

\begin{tikzcd}[column sep=small, row sep=small, nodes={font=\scriptsize}]
\prod_{i} \operatorname{Hom}_{\mathrm{top}}(X, E'_{i}) \arrow[r, hook, "\alpha"'] & \prod_{i} \operatorname{Hom}_{\mathrm{top}}(X, E_{i}) \\
& \operatorname{Hom}_{\mathrm{top}}(X, E) \arrow[u, "\beta"']
\end{tikzcd}
\note{la flèche horizontale porte « pl. fid. », la verticale un « $\simeq$ »}

\page{8}
L'image \add{essentielle} \struck{par} $\alpha$ est la sous-catégorie
\add{str.} pleine de $\prod_{i} \operatorname{Hom}_{\mathrm{top}}(X, E_{i})$
formée des systèmes $(f_{i} : X \to E_{i})$ tel que pour tout $i$, $f_{i}$ se
factorise par $E'_{i}$ ; \ill{} \struck{\ill{}} par $\beta$, et correspond
\ill{} la s.-cat. pleine de $\operatorname{Hom}_{\mathrm{top}}(X, E)$ formée des
$f : X \to E$ tel que $\forall i$, $\mathrm{pr}_{i} \circ f$ se factorise par
$E'_{i}$, i.e. tel que $\forall i$, $f$ se factorise par $F_{i}$, i.e. tels que
\uncertain{$X$} se factorise par $F$. Donc
\[
\prod_{i} \operatorname{Hom}_{\mathrm{top}}(X, E'_{i}) \;\cong\;
\operatorname{Hom}_{\mathrm{top}}(X, F).
\]
\note{la page écrit $F$ là où le sous-topos intersection a été noté $E'$ à la
page précédente}

\textbf{Corollaire de a) et b)} : Le 2-produit d'une famille de topos existe
tjrs.

Pour avoir la $\varprojlim^{\mathrm{top}}$ d'un topos fibré, il faut
$\mathcal{F} \to I$, $\mathcal{F}_{j} \xrightarrow{\ f^{*}\ } \mathcal{F}_{i}$,
$j \xleftarrow{\ f\ } i$ :
\[
\varprojlim_{i} \operatorname{Hom}_{\mathrm{top}}(X, \mathcal{F}_{i})
\;\simeq\;
\varprojlim_{i} \operatorname{Hom}'(\mathcal{F}_{i}, X)
\;\simeq\; \ill{}
\]

$(E_{i})_{i \in I}$ pseudo \struck{syst.} projectif de topos,
$\mathcal{F}_{1} \rightrightarrows \mathcal{F}_{2} \rightrightarrows
\mathcal{F}_{3}$ :
\[
\varprojlim_{i}^{\mathrm{top}} E_{i} \longrightarrow
\prod_{i \in \operatorname{Ob} I} E_{i} \rightrightarrows
\prod_{\alpha \in \mathrm{Fl}_{1} I} E_{b(\alpha)} \rightrightarrows
\prod_{(\alpha,\beta) \in \mathrm{Fl}_{2}(I)} E_{b(\beta)}
\]
les flèches étant $\mathrm{pr}_{s(\alpha)}$, $\mathrm{pr}_{s(\beta)}$ et
$E_{s(\alpha)} \to E_{b(\alpha)}$.
\note{tout ce bas de page est barré de deux longues diagonales ; le troisième
terme porte trois flèches et non deux}

\page{9}
$X$, $Y$ deux espaces topologiques, considérons $\operatorname{Top}(X)
\overset{2}{\times}_{\mathrm{top}} \operatorname{Top}(Y)$ et
$\operatorname{Top}(X \times Y)$ (on vérifie facilement que ce dernier ne
dépend que de $\operatorname{Top}(X)$ et de $\operatorname{Top}(Y)$, i.e. des
\struck{\ill{}} espaces \uncertain{sobres} associés à $X$ et $Y$) d'où un
morphisme canonique de topos
\[
\struck{\text{Prop.}} \quad (*) \qquad
\boxed{\;\operatorname{Top}(X \times Y) \longrightarrow \operatorname{Top}(X)
\overset{2}{\times}_{\mathrm{top}} \operatorname{Top}(Y)\;}
\]
et on veut des conditions pour que ce soit une équivalence.

\marginal{$\mathcal{O}_{X} =$ cat. des ouverts de $X$}

Or on voit que $\operatorname{Top}(X \times Y) = \widetilde{\mathcal{O}_{X}
\times \mathcal{O}_{Y}}$, \struck{où} la topologie à mettre sur
$\mathcal{O}_{X} \times \mathcal{O}_{Y}$ est la « topologie produit » $\pi$,
i.e. celle pour laquelle les faisceaux sont les foncteurs contravariants
$(\mathcal{O}_{X} \overset{2}{\times} \mathcal{O}_{Y})^{\circ} \to
(\mathrm{Ens})$ qui sont des faisceaux \emph{séparément}. \struck{Il} Il y a
aussi la topologie $\pi'$ \struck{\ill{}} « induite » par la précédente
$X \times Y$ en utilisant l'injection (supposons $X \neq \emptyset$ et
$Y \neq \emptyset$) $\mathcal{O}_{X} \times \mathcal{O}_{Y} \subset
\mathcal{O}_{X \times Y}$. Il est clair que tout faisceau pour cette topologie
(faisceau sur $X \times Y$) est un bifaisceau \struck{\ill{}} donc la 2e top.
$\pi'$ est plus fine que $\pi$. On voit déjà que \struck{\ill{}} le
morphisme (*) est un \emph{plongement de topos}, et \struck{\ill{}} est ramené
à voir si ce plongement est une équivalence, i.e. $\pi = \pi'$. Donc on est
ramené à voir si tout recouvrement pour $\pi'$ $U \times V$
($U \in \mathcal{O}_{X}$, $V \in \mathcal{O}_{Y}$) par des $U_{i} \times V_{i}$
[i.e. rec. ensembliste dans $X \times Y$] est un rec. pour $\pi$, i.e.
\struck{si} si
\[
\bigcup_{i} U_{i} \times V_{i} = U \times V \implies
\operatorname*{Sup}_{i} U_{i} \times V_{i} = U \times V
\]
le Sup étant pris dans l'\uncertain{ensemble} des sous-objets de
$U \times V$ \ill{} le topos $\widetilde{\mathcal{O}_{X} \times
\mathcal{O}_{Y}}$ (top. $\pi$).

\textbf{Proposition} Pour qu'il en soit ainsi, il suffit que tout
\struck{\ill{}} point de $U$ (ou de $V$) admette un voisinage (\ill{})
quasi-compact

\textbf{Corollaire} \struck{Si} tous points de $X$ (ou de $Y$) admettant un
syst. fond. de voisinages qui sont quasi-compacts, alors (*) est une
équivalence.

\marginal{NB les $U \times V$ ($U \in \mathcal{O}_{X}$, $V \in
\mathcal{O}_{Y}$) forment une base de la top. de $X \times Y$}

\marginal{NB comme $\operatorname{Top}(X \times Y)$ \ill{}
$\operatorname{Top}(X) \overset{2}{\times}_{\mathrm{top}}
\operatorname{Top}(Y)$ \ill{}}

\note{une longue note marginale, écrite en diagonale le long de toute la marge
de gauche et recoupant les deux précédentes, reste en grande partie illisible}

\page{10}
\textbf{Dém.} Pour tout $x \in X$, soit $U_{x}$ un voisinage
\uncertain{quasi-compact}. Donc $\forall\, y \in Y$ il existe
$I_{x,y} \subset I$, $I_{x,y}$ partie finie, telle que \struck{\ill{}}
$U_{x} \times \{y\}$ soit recouvert par les $U_{i} \times V_{i}$
($i \in I_{x,y}$). Soit $V_{x,y}$ l'intersection \struck{\ill{}} des $V_{i}$,
$V_{x,y} = \bigcap_{i \in I_{x,y}} V_{i}$, donc on a
\[
\begin{cases}
\mathring{U}_{x} = U'_{x} \subset \bigcup_{i \in I_{x,y}} U_{i}, \\
V_{i} \supset V_{x,y} \text{ pour } i \in I_{x,y}
\end{cases}
\]
Donc
\[
S \overset{\text{déf}}{=} \operatorname*{Sup}_{i \in I}(U_{i} \times V_{i})
\;\geqslant\; \operatorname*{Sup}_{i \in I_{x,y}}(U_{i} \times V_{i})
\;\geqslant\; \Bigl(\operatorname*{Sup}_{i \in I_{x,y}} U_{i}\Bigr) \times V_{x,y}
\;\geqslant\; U'_{x} \times V_{x,y}
\]
donc comme ceci est vrai pour tout $y$, on a
\[
S \geqslant \operatorname*{Sup}_{y} \mathring{U}'_{x} \times V_{x,y}
= \mathring{U}'_{x} \times \operatorname*{Sup}_{y} V_{x,y}
= \mathring{U}'_{x} \times V
\]
d'où
\[
S \geqslant \operatorname*{Sup}_{x} (\mathring{U}'_{x} \times V)
= \Bigl(\operatorname*{Sup}_{x} \mathring{U}'_{x}\Bigr) \times V
= U \times V
\]
cqfd.

\textbf{Réciproquement} Supposons que $X$ (ou $Y$) ne contienne pas de parties
\uncertain{loc. fermées} non vides dont aucun point n'ait de voisinage
\add{relatif} quasi-compact. Je dis qu'\struck{alors} a) est une équivalence.

Considérons en effet les ouverts $U'$ dans \struck{$U$} tels que $S \geqslant
U' \times V$ ; il suffit de prouver que $U$ en est parmi. Or la
\struck{\ill{}} \ill{} ouverts, en est un $=$ \uncertain{sup}, donc il y a un
plus grand ouvert \ill{} $=$ \uncertain{sup} [N.B. $\emptyset$ en est un !]

Prouvons $U' = U$. \struck{Soit $Z$} Sinon, la partie loc. fermée $Z$ de
\struck{$U$} \struck{\ill{}} non vide, doit avoir un $x \in Z$ ayant un
voisinage relatif $T_{x} \subset Z$ quasi-compact. Soit $U'' \subset U$ un
ouvert contenant $x$ tel que $U'' \cap Z \subset T$. \struck{Donc} Il existe
$\forall\, y \in Y$ une partie finie $I_{y}$ de $I$ telle que
\[
T \times \{y\} \subset \bigcup_{i \in I_{y}} U_{i} \times V_{i},
\qquad\text{et pour } V_{y} = \bigcap_{i \in I_{y}} V_{i},
\quad \struck{\ill{}}
\]
\struck{$T \times V_{y} \subset \bigcup_{i \in I} U_{i} \times V_{i}$ donc}
$T \subset \bigcup_{i \in I_{y}} U_{i}$, $i \in I_{y} \implies V_{i} \supset
V_{y}$ donc on a
\[
S \geqslant \operatorname*{Sup}_{i \in I_{y}}(U_{i} \times V_{i})
\geqslant \operatorname*{Sup}_{i \in I_{y}} U_{i} \times V_{y}
\geqslant \Bigl(\operatorname*{Sup}_{i \in I_{y}} U_{i}\Bigr) \times V_{y}
\]
Mais on sait déjà que $S \geqslant U' \times V \geqslant U' \times V_{y}$ donc
\[
S \geqslant \operatorname*{Sup}(U'_{i}, U_{i}\ i \in I_{y}) \times V_{y}
\geqslant (U' \cup U'') \times V_{y}
\]

\page{11}
Comme ceci est vrai pour tout $y$, on a
\[
S \geqslant \operatorname*{Sup}_{y}\bigl((U' \cup U'') \times V_{y}\bigr)
= (U' \cup U'') \times \operatorname*{Sup}_{y} V_{y}
= (U' \cup U'') \times V
\]
contredisant au caractère maximal de $U'$, OK.

\subsection*{Cas de produits \struck{infinis} d'une famille quelconque de topos}

On trouve que si tous les facteurs sauf \struck{un seul} \add{au plus} sont
tels que tout pt ait un système fond. de voisinages (\uncertain{q-s vois.
ouverts}) quasi-compacts, et si tous sauf au plus un seul, ont \struck{\ill{}}
voisinages quasi-compacts, alors
\[
\operatorname{Top}\Bigl(\prod_{i \in I} X_{i}\Bigr) \longrightarrow
\prod_{i}^{(2)}{}_{(\mathrm{top})} \operatorname{Top}(X_{i})
\]
est une équivalence de topos.

C'est essentiellement le même argument, travaillant sur la cat. produit
\struck{restreint} $\prod' \mathcal{O}_{X_{i}}$ (familles des
$(U_{i})_{i \in I}$, $U_{i} \in \mathcal{O}_{X_{i}}$ \ill{} $U_{i}$ ;
$U_{i} = X_{i}$ (\struck{\ill{}} pour \add{presque} tout \ill{} de
$\mathcal{O}_{X_{i}}$) pour presque tout $i$) et la « topologie produit » sur ce
produit.

Dans le cas où $I$ est fini, c'est un corollaire immédiat de ce qui précède.
Dans le cas général, on a $I = \varinjlim_{\alpha} J_{\alpha}$ ($J_{\alpha}$
parties finies de $I$) d'où
\begin{align}
\prod_{i}^{\mathrm{top}} \operatorname{Top}(X_{i})
&\cong \varprojlim_{\alpha}{}^{\mathrm{top}}
   \Bigl(\prod_{i \in J_{\alpha}} \operatorname{Top}(X_{i})\Bigr) \nonumber\\
&\cong \varprojlim_{\alpha}{}^{\mathrm{top}}
   \Bigl(\operatorname{Top}\bigl({\textstyle\prod_{i \in J_{\alpha}}} X_{i}\bigr)\Bigr)
   \qquad\Bigl[{\textstyle\prod_{i \in J_{\alpha}}} X_{i} = Z_{\alpha}\Bigr] \nonumber\\
&\approx \Bigl(\underbrace{\varinjlim_{\alpha} \mathcal{O}_{Z_{\alpha}}}_{\mathcal{O}'_{Z}}\Bigr)^{\sim} \nonumber
\end{align}

$\mathcal{O}'_{Z}$ est la catégorie des ouverts de $Z$ \add{$= \prod X_{i}$}
(qui sont images inverses \struck{des} d'ouverts de produits finis
$\prod_{i \in J_{\alpha}} X_{i} = Z_{\alpha}$. Ils la forment \struck{une}
alors la topologie \ill{} (\struck{$\varprojlim$}, qui définit \ill{} le
foncteur précédent, mais aussi la top. plus fine $\pi'$ provenant de la top. de
$Z$. Si tous les $X_{i}$ sauf un nb fini sont quasi-compacts [les autres étant
ce qu'ils veulent] on trouve par un argument analogue \add{à dessus} que les
top. sont

\page{12}
$W \in \mathcal{O}'_{Z}$, donc de la forme $U \times Z'_{\alpha}$, les mêmes.
\struck{Soit en effet} ($U$ un ouvert de $Z_{\alpha}$, et $Z'_{\alpha} =
\prod_{i \in I - J_{\alpha}} X_{i}$, \struck{\ill{}} (d'où $Z = Z_{\alpha}
\times Z'_{\alpha}$), et soient $(V_{\lambda})_{\lambda \in \Lambda}$,
$V_{\lambda} \in \mathcal{O}'_{Z}$, avec $W \subset \bigcup_{\lambda}
V_{\lambda}$ \ill{} prouvons qu'il \struck{\ill{}}
$\operatorname*{Sup}_{\lambda} W \geqslant U \times Z'_{\alpha}$.
\struck{\ill{}} $J_{\alpha}$ \ill{} pour les $X_{i}$ ($i \in I - J_{\alpha}$)
soient \ill{} quelconques, alors $Z'_{\alpha}$ est quasi-compact.
\note{deux taches d'encre couvrent la fin de deux lignes de ce paragraphe}

$\forall\, z \in \struck{\ill{}} W$, $\{z\} \times Z'_{\alpha}$ est recouvert
par un nb fini de $V_{\lambda}$, disons pour $\lambda \in \Lambda_{z}$, donc
$\exists$ \struck{\ill{}} $J_{\beta(z)} \supset J_{\alpha}$ tel que chaque
$V_{\lambda}$ ($\lambda \in \Lambda_{z}$) de la forme $V'_{\lambda} \times
Z'_{\beta(z)}$, avec $V'_{\lambda} \subset \mathcal{O}_{Z_{\beta(z)}}$.
\struck{\ill{}} Comme $\{z\} \times Z'_{\alpha} \subset \bigcup(V'_{\lambda}
\times Z'_{\beta(z)})$, cela \struck{\ill{}} signifie
\[
\{z\} \times \prod_{i \in J_{\beta(z)} - J_{\alpha}} X_{i} \subset
\bigcup_{\lambda \in \Lambda_{z}} V'_{\lambda} ,
\]
donc il existe un voisinage ouvert $U_{z} \ni z$ dans $Z_{\alpha}$ tel que
$U_{z} \times Z'_{\alpha} \subset \bigcup_{\lambda \in \Lambda_{z}}
V_{\lambda}$. Mais alors \struck{\ill{}} on aura l'inclusion
\[
U_{z} \times \prod_{i \in J_{\beta(z)} - J_{\alpha}} X_{i} \subset
\bigcup_{\lambda \in \Lambda_{z}} V'_{\lambda}
\]
d'où
\[
S \overset{\text{déf}}{=} \operatorname*{Sup}_{\lambda \in \Lambda}(V_{\lambda})
\geqslant \operatorname*{Sup}_{\lambda \in \Lambda_{z}} V_{\lambda}
\geqslant U_{z} \times Z'_{\alpha} .
\]
Comme ceci est vrai pour tout $z$, \struck{\ill{}}
\[
S \geqslant \operatorname*{Sup}_{z}\bigl(U_{z} \times Z'_{\alpha}\bigr)
= \Bigl(\operatorname*{Sup}_{z \in Z} U_{z}\Bigr) \times Z'_{\alpha}
= W \times Z'_{\alpha} ,
\]
cqfd.

\subsection*{Contre-exemples pour des produits infinis}

Prenons $X_{i}$ discret pour tout $i$, on aura
\[
\prod_{i \in I} \operatorname{Top}(X_{i}) =
\Bigl(\underbrace{\varinjlim_{\alpha} \mathcal{O}_{Z_{\alpha}}}_{\mathcal{O}'_{Z}}\Bigr)^{\sim} ,
\]
munie par la top.

\page{14}
$\pi$ de $\mathcal{O}'_{Z}$ ainsi que la topologie déduite de la top. produit
sur $Z = \struck{\prod X_{i}} = \varprojlim_{\alpha} Z_{\alpha}$ (si $I$ infini
et les $X_{i}$ ne sont pas presque tous finis). \struck{\ill{}} On a une
famille fond. de voisinages \struck{\ill{}} de $Z$ par un objet de
$\mathcal{O}'$ en associant à tt $z = (x_{i})_{i \in I} \in Z$ une partie finie
$J(z)$ de $J$, \struck{telle} \add{très}
\[
\struck{\ill{}\ W(z) =} \qquad
V_{z} = \{\, z' \in Z \mid z'_{i} = z_{i} \text{ si } i \in J - J(z) \,\} .
\]
\struck{La question est de voir si tous ces $(V_{z})_{z \in Z}$} \add{pas} rec.
pour $\pi$, ou encore
\note{la page s'arrête ici ; le reste du feuillet est vierge, et un « ? » est
porté en marge}

\section*{$\varprojlim$ filtrantes de topos essentiels $\widehat{C}$}

\page{15}
On est intéressé aux $\varprojlim$ pour un topos fibré, i.e. les fibres sont
des topos $\widehat{C}$ ($C$ \struck{\ill{}} \add{arbitr.}) et les morphismes
de transition \struck{$f : \widehat{C}' \to \widehat{C}$} \struck{sont}
essentiels, i.e. associés à des foncteurs
\[
f_{0} : C \longrightarrow C'
\]
Il est \struck{\ill{}} \struck{que alors} $\mathcal{X} =
\varprojlim^{\mathrm{top}} \widehat{C_{i}}$ soit un topos essentiel et ses
projections $\mathcal{X} \to \widehat{C_{i}}$ des morphismes essentiels.

Mais pour le calcul commode de $\mathcal{X}$ on va supposer que les
$\widehat{f}_{ij} : \widehat{C_{j}} \to \widehat{C_{i}}$ \struck{\ill{}}
(associés aux $f_{ij} : C_{j} \to C_{i}$) soient ess., i.e. des morph. de sites
donnés par des foncteurs
\[
g_{ji} : C_{i} \longrightarrow C_{j}
\]
(i.e. induits par $\widehat{f}_{ij}^{*}$ i.e. précis, et adjoints à
\emph{droite} des $f_{ij}$).

On trouve alors
\[
\mathcal{X} = \varprojlim_{i, \widehat{f}_{ij}}{}^{\mathrm{top}} \widehat{C_{i}}
\;\simeq\; \Bigl(\operatorname*{colim}_{i, g_{ij}} C_{i}\Bigr)^{\wedge}
\]
et les morphismes $\mathcal{X} \xrightarrow{\ \mathrm{pr}_{i}\ }
\widehat{C_{i}}$ sont associés aux morphismes de sites donnés par le foncteur
can.
\[
C_{i} \xrightarrow{\ \alpha_{i}\ } C = \operatorname*{colim} C_{j} ,
\]
donc pour $F \in \widehat{C} = \mathcal{X}$, \ill{}
\[
\mathrm{pr}_{i*}(F) = F \circ \alpha_{i} ,
\qquad\text{et pour } G \in \widehat{C_{i}} \ \ill{}
\]
\[
\mathrm{pr}_{i}^{*}(G) =
\]
\note{la formule reste inachevée ; la page s'arrête ici}

\section*{Foncteurs sur un topos $\widetilde{C}$ commutant aux $\varinjlim$}

\page{17}
Soient $C$ un $U$-site, $E$ une catégorie aux petites $\varinjlim$,
considérons les catégories
\[
\operatorname{Hom}'(\widetilde{C}, E)
\qquad
\operatorname{Hom}(C, E)
\]
et soit $\operatorname{Hom}'(\widetilde{C},E)$ \struck{désigner} la
sous-catégorie pleine de $\operatorname{Hom}$ formée des foncteurs qui
commutent aux petites $\varinjlim$. On va définir des foncteurs

\begin{tikzcd}
\operatorname{Hom}(\widetilde{C}, E) \arrow[r, bend left=20, "\alpha"] & \operatorname{Hom}(C,E) \arrow[l, bend left=20, "\beta"]
\end{tikzcd}

en posant
\begin{align}
\alpha(u) &= u \circ \varepsilon_{C} \nonumber\\
\beta(u_{0}) &= \Bigl[F \mapsto \varinjlim_{C/F} u_{0}(X)\Bigr] \nonumber
\end{align}
($\varepsilon_{C} : C \to \widetilde{C}$ le foncteur canonique ; $=
\widehat{\beta} \mid \widetilde{C}$, i.e. $\widehat{\beta} : \widehat{C} \to E$
est le prolongement can. de $\beta$ défini par commutation aux $\varinjlim$)

Soit $u \in \operatorname{Ob} \operatorname{Hom}(\widetilde{C},E)$, on a un hom
\[
i_{u} : \beta\alpha(u) \longrightarrow u
\]
donné par $\forall\, F \in \operatorname{Ob} \widetilde{C}$ par
\[
i_{u}(F) : \bigl(\beta(\alpha(u))\bigr)(F)
= \varinjlim_{C/F} u(\varepsilon_{C}(X)) \longrightarrow u(F)
\]
\marginal{$C \xrightarrow{j_{0}} \widehat{C}$, $\widetilde{C}
\xrightarrow{i_{0}} \widehat{C}$, $a : \widehat{C} \to \widetilde{C}$ ;
$\varepsilon_{C} = a j_{0}$}
\marginal{déduit de l'isom. can. $\varinjlim_{C/F} (\varepsilon_{C}) X
\xrightarrow{\sim} F$ en appliquant $u$}

On trouve que l'homom. $i_{u}(F)$ est bien fonctoriel en $F$, d'où $i_{u}$, et
que $i_{u}$ est fonctoriel en $u$, d'où
\[
\boxed{\;\beta\alpha \xrightarrow{\ i\ } \mathrm{id}_{\operatorname{Hom}(\widetilde{C},E)}\;}
\]

Notons tout de suite que \struck{\ill{}} $i_{u} : \beta\alpha(u) \to u$ est un
isom. sss $u : \widetilde{C} \to E$ commute aux $\varinjlim$ du type
$\varinjlim_{C/F} X \xrightarrow{\sim} F$ ($F \in \operatorname{Ob}
\widetilde{C}$), \struck{\ill{}} \add{donc que} $u \in \operatorname{Ob}
\operatorname{Hom}'(\widetilde{C}, E)$.

Soit $u_{0} \in \operatorname{Hom}(C,E)$, on va définir
\[
j_{u_{0}} : \struck{\ill{}}\; u_{0} \longrightarrow \alpha\beta(u_{0})
\]
\struck{\ill{}} \add{donné} pour $Z \in \operatorname{Ob} C$, par
\[
j_{u_{0}}(Z) : \quad u_{0}(Z) \longrightarrow
\bigl(\alpha(\beta(u_{0}))\bigr)(Z)
= \beta(u_{0})(\varepsilon_{C}(Z))
= \varinjlim_{C/\varepsilon_{C}(Z)} u_{0}(X)
\]
en notant que $(Z, \text{morphisme can. de } Z \text{ dans }
\varepsilon_{C}(Z))$ est un objet de $C/\varepsilon_{C}(Z)$, soit $X_{0}$, et
par suite $u_{0}(X_{0}) = u_{0}(Z)$ s'envoie dans la $\varinjlim$.

\page{18}
On vérifie que cette application est fonctorielle en $Z$, donc définit
$j_{u_{0}} : u_{0} \to \alpha\beta(u_{0})$, et que cette dernière est
fonctorielle en $u_{0}$, d'où un hom. de foncteurs
\[
\boxed{\;j : \mathrm{id}_{\operatorname{Hom}(C,E)} \longrightarrow \alpha\beta\;}
\]
\struck{Signalons} Notons tout de suite que pour \struck{que} ce dernier,
$j_{u_{0}} : u_{0} \to \alpha\beta(u_{0})$, soit un isom. sss
\struck{\ill{}} $\forall\, Z \in \operatorname{Ob} C$,
\[
u_{0}(Z) \longrightarrow \varinjlim_{C/\varepsilon_{C}(Z)} u_{0}(X)
\]
est un isom, ou encore sss pour tout tel $Z$ et $\xi \in \operatorname{Ob} E$,
la flèche déduite de la précédente en appliquant $\operatorname{Hom}(-, \xi)$
est bijective, i.e.
\[
\varprojlim_{C/\varepsilon_{C}(Z)} \operatorname{Hom}(u_{0}(X), \xi)
\longrightarrow \operatorname{Hom}(u_{0}(Z), \xi)
\quad\text{bijective.}
\]
On sait \struck{\ill{}},
\[
u'_{0} : E \longrightarrow \widehat{C}
\]
le foncteur défini par
\[
u'_{0}(\xi)(Y) = \operatorname{Hom}(u_{0}(Y), \xi)
\]
de sorte que l'application précédente devient, posant $P = u'_{0}(\xi)$ :
\[
\varprojlim_{C/\varepsilon_{C}(Z)} P(X) \longrightarrow P(Z)
\qquad\text{un isom.}
\]
\struck{\ill{}} Cette condition est vide si $\varepsilon_{C}$ est pl. fidèle
i.e. la top. de $C$ est moins fine que la top. canonique (car alors
$C/\varepsilon_{C}(Z) \simeq C/Z$ admet $Z$ comme objet final).

On va prouver que $i$, $j$ sont des homom. d'adjonction associés pour les
foncteurs $\alpha$ et $\beta$. [Ce qui prouve que $\alpha$ et $\beta$
\struck{induisent} \struck{établissent} des équivalences entre les
sous-catégories \add{str. pleines} de $\operatorname{Hom}(\widetilde{C},E)$ et
$\operatorname{Hom}(C,E)$ qu'on a décrites au préal. On \struck{\ill{}}
\struck{\ill{}} (on a induit un foncteur pl. fidèle
\[
\operatorname{Hom}'(\widetilde{C}, E) \longrightarrow \operatorname{Hom}(C,E)
\]
dont l'image essentielle est formée des \struck{\ill{}} $u_{0} : C \to E$
satisfaisant la condition ci-dessus, et telle que $\beta(u_{0}) \in
\operatorname{Hom}'(\widetilde{C}, E)$ — condition qu'on remplacera par une
autre plus simple]

\page{19}
$C$ un $U$-site, $E$ catégorie où les petites $\varinjlim$ sont repr.
Déterminer les foncteurs
\[
u : \widetilde{C} \longrightarrow E
\]
commutant aux $\varinjlim$ (petites) ? Comme \ill{}
\[
u \circ \varepsilon_{C} = u_{0} , \qquad u_{0} : C \longrightarrow E
\]
\marginal{$C \xrightarrow{\varepsilon_{C}} \widetilde{C} \xrightarrow{u} E$, le
composé étant $u_{0}$}
car pour $F \in \widetilde{C}$, on a un isom. fonctoriel en $F$
\[
F = \varinjlim_{C/F} \varepsilon_{C}(X)
\]
d'où
\[
u(F) = \varinjlim_{C/F} u_{0}(X) .
\]
De façon précise, on trouve un foncteur $u \mapsto u \circ \varepsilon_{C}$
\[
a : \operatorname{Hom}'(\widetilde{C}, E) \longrightarrow \operatorname{Hom}(C,E)
\]
qui est pl. fid., le signe $\operatorname{Hom}'$ désignant les foncteurs
commutant aux petites $\varinjlim$. En fait, on a un foncteur
« prolongement à $\widetilde{C}$ »
\[
\beta : \operatorname{Hom}(C,E) \longrightarrow \operatorname{Hom}(C,E)
\]
\note{le but devrait être $\operatorname{Hom}(\widetilde{C},E)$ ; la page écrit
$C$ des deux côtés}
défini par
\[
\beta(u_{0}) = F \mapsto \varinjlim_{C/F} u_{0}(X)
\]
et on a un isom. canonique
\[
\beta \circ \alpha \simeq i
\qquad (\text{inclusion de } \operatorname{Hom}'(\widetilde{C},E)
\hookrightarrow \operatorname{Hom}(\widetilde{C},E))
\]
Mais il n'est pas clair que $\forall\, u_{0} : C \to E$, $\beta(u_{0})$ soit
dans $\operatorname{Hom}'$. À vrai dire, $\beta(u_{0})$ est induit par
$\widehat{\beta}(u_{0}) : \widehat{C} \to E$ défini de la $=$ façon, et on sait
que $\widehat{\beta}(u_{0})$ commute aux $\varinjlim$, on a donc

\begin{tikzcd}
\operatorname{Hom}(C,E) \arrow[r, bend left=20, "\widehat{\beta}"] & \operatorname{Hom}'(\widehat{C}, E) \arrow[l, bend left=20, "\widehat{\alpha}"]
\end{tikzcd}

(cf. rappel plus bas). Identifions, dorénavant (une double équivalence)
$\operatorname{Hom}(C,E)$ à $\operatorname{Hom}'(\widehat{C},E)$, alors
$\alpha$ et $\beta$ deviennent des foncteurs $\alpha'$, $\beta'$
\begin{align}
\operatorname{Hom}'(\widetilde{C}, E) &\xrightarrow{\ \alpha'\ } \operatorname{Hom}'(\widehat{C}, E) \nonumber\\
\operatorname{Hom}'(\widehat{C}, E) &\xrightarrow{\ \beta'\ } \operatorname{Hom}(\widetilde{C}, E) \nonumber
\end{align}

\page{20}
qui s'explicitent ainsi : soit ($a : \widehat{C} \to \widetilde{C}$ le foncteur
« faisceau associé »), donc $\varepsilon_{C} = a\, i_{C}$ ($i_{C} : C \to
\widehat{C}$) donc pour $u \in \operatorname{Hom}'(\widetilde{C},E)$,
$u \circ \varepsilon_{C} = (ua)\, i_{C}$, donc (\struck{\ill{}} que $a$, donc
$ua$, commute aux petites $\varinjlim$)
\[
\alpha'(u) = u \circ a
\]
\struck{Donc} D'autre part, on \ill{} que pour $v \in
\operatorname{Hom}'(\widehat{C}, E)$ (correspondant à $u_{0} : C \to E$)
\[
\beta'(v) = v \circ j_{a} = v \mid \widetilde{C}
\qquad (\text{où } j : \widetilde{C} \to \widehat{C} \text{ est l'inclusion})
\]
On a bien
\[
\beta'\alpha'(u) = u \circ a \mid \widetilde{C} \simeq u
\]
Le fait que $\alpha'$ soit pl. fidèle provient du fait que $\widetilde{C}$ est,
moyennant $a$, une catégorie de quotients de $\widehat{C}$, de sorte que
$\operatorname{Hom}(\widetilde{C},E)$ (sans $'$) s'identifie à la
sous-catégorie pleine de $\operatorname{Hom}(\widehat{C},E)$ formée des
foncteurs $F$ qui sont tels que
\[
\struck{\ill{}}\quad u \in \mathrm{Fl}\,\widehat{C},\ a(u) \text{ inversible}
\implies F(u) \text{ inversible.}
\]
Est-il vrai que dans cette identification, les foncteurs qui commutent aux
$\varinjlim$ petites se correspondent ? Il vaut à voir que si $\mathcal{U} :
\widetilde{C} \to E$ est tel que $\mathcal{U} \circ a$ commute aux
$\varinjlim$, \ill{} on aura $\mathcal{U} = (\mathcal{U} \circ a)\, i$, et si
\ill{} $X_{\alpha}$ de $\widetilde{C}$,
$\varinjlim_{\alpha}^{\widetilde{C}} X_{\alpha} =
a\bigl(\varinjlim_{\alpha}^{\widehat{C}} X_{\alpha}\bigr)$, donc
\[
\mathcal{U}\Bigl(\varinjlim_{\alpha}{}^{\widetilde{C}} X_{\alpha}\Bigr)
= \mathcal{U}\, a \Bigl(\varinjlim_{\alpha}{}^{\widehat{C}} X_{\alpha}\Bigr)
= (\mathcal{U}a)\Bigl(\varinjlim_{\alpha} X_{\alpha}\Bigr)
= \varinjlim_{\alpha} (\mathcal{U} \circ a)(X_{\alpha})
= \varinjlim_{\alpha} \mathcal{U}(X_{\alpha})
\]
O.K.

Donc on trouve que $\alpha'$ induit
\[
\operatorname{Hom}'(\widetilde{C}, E) \xrightarrow{\ \simeq\ }
\text{sous-catégorie pleine de } \operatorname{Hom}(\widehat{C}, E)
\]
formée des foncteurs qui a) commutent aux $\varinjlim$ petites, b)
transforment les morphismes « bicouvrants » de $\widehat{C}$ (i.e. tels que
$a(u)$ inv.) en isom.

Donc on cherche
\[
\operatorname{Hom}'(\widetilde{C}, E) \xrightarrow[\ \simeq\ ]{\ \alpha_{0}\ }
\struck{\ill{}} \text{ sous-catégorie pleine de } \operatorname{Hom}(C, \struck{\widetilde{E}}\,E)
\]
formée des $u_{0} : C \to E$ tels que $\widehat{u}_{0} : \widehat{C} \to E$
transforme les morphismes « bicouvrants » en isom.

Mais cette condition implique que tt famille couvrante $X_{\alpha} \to X$ de
$\widetilde{C}$ (a fortiori d'une \struck{famille} famille épimorphique de
$\widetilde{C}$) soit transformée en famille épimorphique de $E$. Cela
est-il suffisant ? Supposons cette condition remplie,

\marginal{non}
\note{ce « non » est porté au crayon en bas de page, et répond à la question
qui précède}

\end{document}
